% ═══════════════════════════════════════════════════════════════
% LERNBLATT - Physik Klasse 6: Bewegung
% Summarisch (ausgefüllt) - als Unterlage für Stundenleistung
% ═══════════════════════════════════════════════════════════════

\documentclass[11pt, a4paper]{article}

\usepackage[utf8]{inputenc}
\usepackage[T1]{fontenc}
\usepackage[ngerman]{babel}
\usepackage{helvet}
\renewcommand{\familydefault}{\sfdefault}

\usepackage[margin=1.5cm, top=1.8cm, bottom=1.5cm]{geometry}
\usepackage{enumitem}
\usepackage{tabularx}
\usepackage{array}
\usepackage{multicol}
\usepackage{tikz}
\usepackage{amssymb}
\usepackage{amsmath}
\usepackage{fancyhdr}
\usepackage{xcolor}
\usepackage{tcolorbox}
\tcbuselibrary{skins, breakable}

% Farben
\definecolor{physik}{HTML}{2c5aa0}
\definecolor{physik-hell}{HTML}{e8f0fa}
\definecolor{hellgrau}{HTML}{F5F5F5}

% Kopfzeile
\pagestyle{fancy}
\fancyhf{}
\fancyhead[L]{\small\textcolor{physik}{\textbf{Physik Klasse 6}} | Lernblatt Bewegung}
\fancyhead[R]{\small Name: \underline{\hspace{4cm}}}
\fancyfoot[C]{\small Seite \thepage}
\renewcommand{\headrulewidth}{0pt}

% Boxen
\newtcolorbox{merkkasten}[1][]{
    colback=physik-hell,
    colframe=physik,
    boxrule=1.5pt,
    arc=0pt,
    left=8pt, right=8pt, top=8pt, bottom=6pt,
    title={\small\textbf{#1}},
    coltitle=white,
    fonttitle=\sffamily,
    attach boxed title to top left={yshift=-2mm, xshift=5mm},
    boxed title style={colback=physik, arc=0pt}
}

\newtcolorbox{formelbox}{
    colback=white,
    colframe=physik,
    boxrule=2pt,
    arc=0pt,
    left=10pt, right=10pt, top=10pt, bottom=10pt
}

% Befehle
\newcommand{\seitentitel}[1]{%
    \noindent\textcolor{physik}{\rule{\textwidth}{2pt}}\\[0.3em]
    {\LARGE\bfseries\textcolor{physik}{#1}}\\[0.2em]
    \textcolor{physik}{\rule{\textwidth}{0.5pt}}
    \vspace{0.5em}
}

\newcommand{\abschnitt}[1]{%
    \vspace{6pt}%
    \noindent{\large\bfseries\textcolor{physik}{#1}}%
    \vspace{4pt}%
}

\setlength{\parindent}{0pt}
\setlength{\parskip}{3pt}

\begin{document}

% ═══════════════════════════════════════════════════════════════
% SEITE 1: BEWEGUNGSARTEN
% ═══════════════════════════════════════════════════════════════

\seitentitel{Teil 1: Bewegungsarten}

\begin{merkkasten}[Kasten 1: Was ist Bewegung?]
Ein Körper ist in \textbf{Bewegung}, wenn er seinen \textbf{Ort} verändert.

Wir beschreiben Bewegungen mit zwei Eigenschaften:
\begin{itemize}[leftmargin=1.5em, itemsep=1pt]
    \item Die \textbf{Bahnform} beschreibt, wie der \textbf{Weg} aussieht.
    \item Die \textbf{Bewegungsart} beschreibt, ob sich die \textbf{Geschwindigkeit} ändert.
\end{itemize}
\end{merkkasten}

\vspace{0.3em}

\begin{minipage}[t]{0.48\textwidth}
\begin{merkkasten}[Kasten 2: Bahnformen]
\textbf{Geradlinig:}\\
Der Körper bewegt sich auf einer \textbf{geraden Linie}.

\begin{center}
\begin{tikzpicture}[scale=0.6]
    \draw[thick, physik, ->] (0,0) -- (3,0);
    \node[below, font=\scriptsize] at (1.5,-0.2) {z.B. Aufzug, Rolltreppe};
\end{tikzpicture}
\end{center}

\textbf{Krummlinig:}\\
Der Körper bewegt sich auf einer \textbf{gebogenen Bahn}.

\begin{center}
\begin{tikzpicture}[scale=0.6]
    \draw[thick, physik, ->] (0,0) arc (180:0:0.8);
    \node[below, font=\scriptsize] at (0.8,-0.2) {z.B. Karussell, Schaukel};
\end{tikzpicture}
\end{center}

\textit{Kreisförmig ist ein Sonderfall von krummlinig!}
\end{merkkasten}
\end{minipage}
\hfill
\begin{minipage}[t]{0.48\textwidth}
\begin{merkkasten}[Kasten 3: Bewegungsarten]
\textbf{Gleichförmig:}\\
Die Geschwindigkeit bleibt \textbf{gleich}.

\begin{center}
\begin{tikzpicture}[scale=0.5]
    \foreach \x in {0,1,2,3} { \fill[physik] (\x,0) circle (4pt); }
    \node[below, font=\scriptsize] at (1.5,-0.3) {gleiche Abstände};
\end{tikzpicture}
\end{center}

\textbf{Ungleichförmig:}\\
Die Geschwindigkeit \textbf{ändert} sich.

\begin{center}
\begin{tikzpicture}[scale=0.5]
    \fill[physik] (0,0) circle (4pt);
    \fill[physik] (0.3,0) circle (4pt);
    \fill[physik] (0.8,0) circle (4pt);
    \fill[physik] (1.5,0) circle (4pt);
    \fill[physik] (2.5,0) circle (4pt);
    \node[below, font=\scriptsize] at (1.25,-0.3) {verschiedene Abstände};
\end{tikzpicture}
\end{center}

\textit{Beispiele: Anfahren, Bremsen}
\end{merkkasten}
\end{minipage}

\vspace{0.5em}

\abschnitt{Beispiele aus dem Alltag}

\begin{center}
\renewcommand{\arraystretch}{1.3}
\begin{tabular}{|p{4.2cm}|c|c|}
\hline
\textbf{Beispiel} & \textbf{Bahnform} & \textbf{Bewegungsart} \\
\hline
Zug auf gerader Strecke & geradlinig & gleichförmig \\
\hline
Auto beim Bremsen & geradlinig & ungleichförmig \\
\hline
Gondel am Riesenrad & krummlinig & gleichförmig \\
\hline
Achterbahn im Looping & krummlinig & ungleichförmig \\
\hline
Sekundenzeiger einer Uhr & krummlinig & gleichförmig \\
\hline
Fallender Stein & geradlinig & ungleichförmig \\
\hline
\end{tabular}
\end{center}

% ═══════════════════════════════════════════════════════════════
% SEITE 2: GESCHWINDIGKEIT
% ═══════════════════════════════════════════════════════════════

\newpage

\seitentitel{Teil 2: Geschwindigkeit}

\begin{merkkasten}[Kasten 4: Die Geschwindigkeit]
Die \textbf{Geschwindigkeit} $v$ gibt an, welchen \textbf{Weg} $s$ ein Körper in einer bestimmten \textbf{Zeit} $t$ zurücklegt.
\end{merkkasten}

\vspace{0.3em}

\begin{formelbox}
\begin{center}
{\Large $v = \dfrac{s}{t}$}

\vspace{0.8em}

\renewcommand{\arraystretch}{1.4}
\begin{tabular}{|c|c|c|}
\hline
\textbf{Größe} & \textbf{Formelzeichen} & \textbf{Einheit} \\
\hline
Geschwindigkeit & $v$ & m/s (Meter pro Sekunde) \\
\hline
Weg (Strecke) & $s$ & m (Meter) \\
\hline
Zeit & $t$ & s (Sekunde) \\
\hline
\end{tabular}
\end{center}
\end{formelbox}

\vspace{0.5em}

\begin{merkkasten}[Kasten 5: Formel umstellen]
\begin{center}
\begin{tabular}{ccc}
\textbf{Gesucht: $v$} & \textbf{Gesucht: $s$} & \textbf{Gesucht: $t$} \\[0.5em]
$v = \dfrac{s}{t}$ & $s = v \cdot t$ & $t = \dfrac{s}{v}$ \\
\end{tabular}
\end{center}
\end{merkkasten}

\vspace{0.5em}

\abschnitt{Beispielrechnungen}

\begin{minipage}[t]{0.48\textwidth}
\textbf{Beispiel 1: Geschwindigkeit berechnen}

Ein Läufer legt 100\,m in 10\,s zurück.

\textbf{Gegeben:} $s = 100$\,m, $t = 10$\,s\\
\textbf{Gesucht:} $v$\\
\textbf{Formel:} $v = \dfrac{s}{t}$\\
\textbf{Einsetzen:} $v = \dfrac{100\,\text{m}}{10\,\text{s}}$\\
\textbf{Ergebnis:} $v = 10$\,m/s
\end{minipage}
\hfill
\begin{minipage}[t]{0.48\textwidth}
\textbf{Beispiel 2: Weg berechnen}

Ein Radfahrer fährt mit 5\,m/s für 20\,s.

\textbf{Gegeben:} $v = 5$\,m/s, $t = 20$\,s\\
\textbf{Gesucht:} $s$\\
\textbf{Formel:} $s = v \cdot t$\\
\textbf{Einsetzen:} $s = 5\,\text{m/s} \cdot 20\,\text{s}$\\
\textbf{Ergebnis:} $s = 100$\,m
\end{minipage}

\vspace{0.8em}

\abschnitt{Typische Geschwindigkeiten}

\begin{center}
\renewcommand{\arraystretch}{1.3}
\begin{tabular}{|l|c|l|}
\hline
\textbf{Was?} & \textbf{Geschwindigkeit} & \textbf{Bedeutung} \\
\hline
Schnecke & 0,001 m/s & 1 mm pro Sekunde \\
\hline
Fußgänger & 1--2 m/s & ca. 1,5 Meter pro Sekunde \\
\hline
Jogger & 3--4 m/s & ca. 3 Meter pro Sekunde \\
\hline
Sprinter (Usain Bolt) & 10 m/s & 10 Meter pro Sekunde \\
\hline
Radfahrer & 5--8 m/s & ca. 6 Meter pro Sekunde \\
\hline
Auto in der Stadt & 14 m/s & ca. 50 km/h \\
\hline
\end{tabular}
\end{center}

\vspace{0.5em}

\begin{merkkasten}[Merke: So gehst du vor!]
\begin{enumerate}[leftmargin=1.5em, itemsep=2pt]
    \item \textbf{Gegeben} aufschreiben (mit Einheiten!)
    \item \textbf{Gesucht} notieren
    \item \textbf{Formel} hinschreiben
    \item \textbf{Einsetzen} (Zahlen in die Formel)
    \item \textbf{Ergebnis} mit Einheit aufschreiben
\end{enumerate}
\end{merkkasten}

\end{document}
